% Chapter 2: Literature Review

\section{Introduction}

This chapter reviews existing literature relevant to the Lesaka Multi-Agent Orchestration System. The review focuses on four key areas: multi-agent systems, graph-based orchestration algorithms, agent specialization, and self-healing mechanisms. The review identifies gaps in current research and establishes the theoretical foundation for the project.

\section{Multi-Agent Systems}

\subsection{Agent Architectures}

Multi-agent systems (MAS) are composed of multiple interacting agents that can perceive, reason, and act autonomously \cite{wooldridge2009introduction}. Agents in a MAS can be homogeneous (similar capabilities) or heterogeneous (different specializations), and can cooperate, coordinate, or negotiate to achieve common or individual goals.

Jennings and Wooldridge \cite{jennings1998agent} identify several key characteristics of MAS:
- Autonomy: Agents operate without direct human intervention
- Social ability: Agents interact with other agents
- Reactivity: Agents respond to environmental changes
- Pro-activeness: Agents take initiative to achieve goals

\subsection{Agent Communication}

Agent communication is essential for coordination in MAS. Common communication paradigms include:
- Direct message passing
- Blackboard systems
- Contract net protocols
- Shared knowledge bases

Research by Sycara \cite{sycara1998multiagent} emphasizes the importance of efficient communication protocols for large-scale MAS, particularly in distributed environments with limited connectivity.

\section{Graph-Based Orchestration}

\subsection{Graph Algorithms in MAS}

Graph-based orchestration uses graph theory to model agent interactions and routing decisions. Nodes represent agents or processing stages, and edges represent communication channels or data flow \cite{diest2010graph}.

Common graph algorithms used in MAS include:
- Shortest path algorithms for routing
- Minimum spanning trees for connectivity
- Maximum flow algorithms for resource allocation
- Graph coloring for conflict resolution

\subsection{Directed Cyclic Graphs (DCG)}

Directed Cyclic Graphs allow for feedback loops and iterative processing, which are essential for self-healing and adaptive behavior in MAS \cite{ghosh2007self}. DCGs enable agents to route data back through previous stages for reprocessing or error recovery.

Research by Ghosh et al. \cite{ghosh2007self} demonstrates that DCG-based orchestration provides superior fault tolerance compared to linear pipelines, particularly in distributed systems with unreliable communication.

\section{Agent Specialization}

\subsection{Domain-Specific Agents}

Specialized agents have expertise in specific domains and can provide more accurate analysis than general-purpose agents. In agricultural contexts, specialized agents can focus on:
- Biomedical analysis (health monitoring)
- Environmental analysis (weather, conditions)
- Market analysis (pricing, timing)
- Resource optimization (feed, water)

Research by Luck et al. \cite{luck2003agent} demonstrates that specialized agents can achieve higher accuracy in domain-specific tasks compared to general-purpose agents.

\subsection{Agent Collaboration}

Specialized agents must collaborate effectively to achieve system-level goals. Collaboration mechanisms include:
- Task allocation
- Result aggregation
- Conflict resolution
- Consensus building

Research by Kraus \cite{kraus1997negotiation} emphasizes the importance of negotiation protocols for effective collaboration among specialized agents.

\section{Self-Healing Mechanisms}

\subsection{Fault Tolerance in MAS}

Self-healing mechanisms enable MAS to recover from failures without human intervention. Key approaches include:
- Redundancy: Multiple agents for critical tasks
- Checkpointing: State recovery from saved states
- Fallback mechanisms: Alternative processing paths
- Adaptive behavior: Dynamic reconfiguration

Research by Ghosh et al. \cite{ghosh2007self} identifies self-healing as a critical requirement for MAS in distributed environments with unreliable communication.

\subsection{Supervisor Nodes}

Supervisor nodes provide centralized oversight and fallback capabilities in MAS. They can:
- Monitor agent health
- Detect failures
- Initiate recovery procedures
- Provide neutral fallback states

Research by Parunak et al. \cite{parunak2001practical} demonstrates that supervisor nodes significantly improve system reliability in industrial MAS applications.

\section{Agricultural Multi-Agent Systems}

\subsection{Livestock Monitoring MAS}

Several MAS have been developed for livestock monitoring. Riquelme et al. \cite{riquelme2018smonitor} describe a MAS for cattle health monitoring using wireless sensor networks. Their system uses multiple agents for data collection, analysis, and alert generation.

However, many existing agricultural MAS lack:
- Graph-based orchestration
- Self-healing capabilities
- Advanced agent specialization
- Integration with data governance frameworks

\subsection{Challenges in Agricultural MAS}

Implementing MAS in agricultural contexts presents unique challenges:
- Limited connectivity in rural areas
- Resource constraints on edge devices
- Environmental interference with communication
- Need for real-time decision-making

Research by Dara et al. \cite{dara2018precision} emphasizes the need for fault-tolerant MAS that can operate reliably in resource-constrained agricultural environments.

\section{Research Gaps}

Based on the literature review, several research gaps have been identified:

\subsection{Gap 1: Graph-Based Orchestration in Agricultural MAS}

While graph algorithms have been extensively studied in general MAS contexts, their application to agricultural IoT is limited. The specific challenges of agricultural telemetry processing (e.g., intermittent connectivity, resource constraints) require tailored graph-based approaches \cite{diest2010graph}.

\subsection{Gap 2: Self-Healing in Resource-Constrained Environments}

Self-healing mechanisms have been studied in industrial MAS, but their application to resource-constrained agricultural environments is limited. The trade-offs between self-healing overhead and resource constraints require further investigation \cite{ghosh2007self}.

\subsection{Gap 3: Agent Specialization for Agricultural Contexts}

Agent specialization has been studied in various domains, but agricultural-specific specialization (e.g., biomedical analysis for cattle, environmental analysis for livestock) requires further research \cite{luck2003agent}.

\subsection{Gap 4: Integration with Data Governance}

Most agricultural MAS focus on data processing and decision-making, but do not integrate with comprehensive data governance frameworks. The integration of MAS with data validation, RBAC, and privacy protection requires further study \cite{riquelme2018smonitor}.

\section{Comparison of Existing Systems}

Table \ref{tab:comparison} compares existing agricultural MAS in terms of algorithmic features.

\begin{table}[H]
\centering
\caption{Comparison of Agricultural Multi-Agent Systems}
\label{tab:comparison}
\begin{tabular}{|l|c|c|c|c|}
\hline
\textbf{System} & \textbf{Graph Routing} & \textbf{Specialized Agents} & \textbf{Self-Healing} & \textbf{State Machine} \\
\hline
Riquelme et al. \cite{riquelme2018smonitor} & No & Partial & No & No \\
Dara et al. \cite{dara2018precision} & No & Yes & Partial | Yes \\
\textbf{Lesaka AI (Proposed)} & \textbf{Yes} & \textbf{Yes} & \textbf{Yes} & \textbf{Yes} \\
\hline
\end{tabular}
\end{table}

The comparison demonstrates that existing systems lack comprehensive algorithmic features. The Lesaka AI system aims to address these gaps by implementing graph-based routing, specialized agents, self-healing mechanisms, and state machine design.

\section{Conclusion}

The literature review has established the theoretical foundation for the Lesaka Multi-Agent Orchestration System. The review has identified key concepts in multi-agent systems, graph algorithms, agent specialization, and self-healing mechanisms. The research gaps identified provide justification for the project's focus on graph-based orchestration, specialized agents, and self-healing for agricultural IoT in Botswana's context.

The next chapter will describe the methodology used to design and implement the Lesaka Multi-Agent Orchestration System.
